\documentclass[12pt]{article}
\usepackage[utf8]{inputenc}
\usepackage[T1]{fontenc}
\usepackage{hyperref}
\title{Makhuwa-Marrevone Language Resources for K-12 Curriculum Development}
\author{Compiled by Claude Sonnet 5.5 (claude-sonnet-5-5) for Andr\'as Kornai}
\date{}
\begin{document}
\maketitle

\section{Introduction}
Makhuwa-Marrevone (Marrevone, Marrevoni; ISO 639-3 xmc; Glottocode makh1262) is a Bantu variety of the coastal Makhuwa cluster in northern Mozambique, in the Moma District area. Population estimates range from about 438,000 to 917,000 depending on the source. The figures may double-count neighbouring varieties. Standard Makhuwa is based on Central Makhuwa, and mutual intelligibility among varieties is limited. This report is short because few Marrevone-specific resources turned up.

\section{Primary Educational Resources}
No Marrevone primers, textbooks or school programmes were found in the search results. The Wikipedia article on Makhuwa describes the wider language and may help locate materials for related varieties.

\section{Linguistic and Reference Materials}
\begin{itemize}
\item Wikipedia, Makhuwa language (varieties, classification, caveats on speaker figures): \url{https://en.wikipedia.org/wiki/Makhuwa_language}
\end{itemize}

\section{Religious and Community Resources}
\begin{itemize}
\item Wycliffe South Africa, Marrevoni (Makhuwa Cluster) translation project, run with the Nampula Roman Catholic Church. Genesis, Luke and Acts are reported complete, with the New Testament as the goal: \url{https://wycliffe.org.za/en/marrevoni}
\item Joshua Project, Makhuwa-Marrevone in Mozambique: \url{https://joshuaproject.net/people_groups/13421}
\end{itemize}

\section{Digital and Audio Resources}
\begin{itemize}
\item Global Recordings Network, Makhuwa-Marrevone: \url{https://globalrecordings.net/en/language/xmc}
\end{itemize}

\section{Curriculum Development Context}
The sources note that most Marrevoni people are Muslim, with a Christian minority, and live by farming and fishing. The translated scripture portions could serve as a source of written text. A curriculum would need an orthography decision, ideally aligned with standard Makhuwa, and community consultation. The search results did not show any existing school use.
\end{document}